\section{Introdução}

As mudanças climáticas têm refletido na incidência de desastres ambientais no Brasil. Apenas no primeiro semestre de 2026, registraram-se 19.462 focos de incêndio \cite{poder360_2026}, complementando os registros de 2024, quando o fogo consumiu 30 milhões de hectares no país, predominantemente sobre florestas nativas na Amazônia \cite{g1_mapbiomas_2025}. A detecção isolada de anomalias térmicas satelitais, contudo, necessita ser correlacionada de forma contínua às vulnerabilidades demográficas e infraestruturais para subsidiar mitigações eficazes.

A integração de bases institucionais (INPE \cite{inpe_bdqueimadas_2026}, FUNAI \cite{funai_terras_2026}, ANEEL \cite{aneel_sigel_2026}, IBGE \cite{ibge_malha_2022}, DATASUS \cite{datasus_cnes_2026}) aos dados de satélite enfrenta profundos desafios de latência e padronização geométrica. Plataformas consolidadas, como o Painel do Fogo \cite{censipam_painel_2026}, prestam suporte essencial, porém arquiteturas de processamento baseadas em lotes (\textit{batch}) apresentam gargalos de vazão ao cruzar milhões de registros geoespaciais e séries temporais de forma interativa. Agregar massas de dados dinâmicas para a geração de respostas analíticas instantâneas permanece um limite técnico em sistemas tradicionais.

Diante dessa barreira, bancos de dados analíticos colunares (OLAP) orientados a \textit{streaming}, como o Apache Druid, despontam como soluções ideais para a absorção paralela de \textit{logs} em tempo real. Contudo, apesar de sua eficiência arquitetural no isolamento entre armazenamento e consulta \cite{yang_druid_2014}, ainda é incerto o impacto na latência quando tais motores são submetidos a volumétricas e complexas junções espaciais interativas (\textit{Lookups}) operando sob restrições severas de hardware e limitação de RAM.

Este trabalho tem como objetivo principal avaliar o desempenho e a estabilidade de processamento do Apache Druid frente a um rigoroso modelo de estresse analítico de dados geoespaciais e ambientais. Para ancorar o experimento em uma realidade aplicada, utiliza-se o SIRDA (Sistema de Inteligência e Resposta a Desastres Ambientais) como estudo de caso, implementando um fluxo \textit{DataOps} contínuo \cite{ereth_dataops_2018} capaz de consolidar riscos topológicos para a Defesa Civil.

O \textit{benchmark} estruturado revela que a otimização na compactação dos dados em segmentos possui um impacto superior ao aumento bruto de poder de processamento, registrando um ganho de 38,8 vezes na vazão de resposta frente a estratégias fragmentadas sob o mesmo hardware. Adicionalmente, os resultados destacam o fenômeno computacional \textit{Too Fast to Cache}, demonstrando que o sistema é capaz de resolver filtragens relacionais em tempos sub-segundo tão rapidamente que a contenção imposta à memória RAM pelo motor interno torna-se um limite arquitetural intransponível antes mesmo que as CPUs atinjam a saturação plena.
